% fig05.tex -- Set A (Architectural) Figure 5: Anchor domain and
% whole-table claim scope.
%
% Cleaner integrated diagram: outer scope(C) rounded rectangle with
% generous interior; ellipse names placed ABOVE each ellipse (outside,
% inside the rectangle); concrete keys inside each ellipse; source-log
% axis below; four numbered conjuncts as circled numerals.

\centering
\begin{tikzpicture}[setA, x=1cm, y=1cm,
    conjNum/.style={
      draw=gDark, line width=0.5pt, circle, fill=gPale,
      minimum size=4.5mm, inner sep=0pt, font=\sffamily\bfseries\scriptsize,
      text=gDark,
    },
    keyDot/.style={
      font=\sffamily\small\bfseries, text=gDark, inner sep=1pt,
    },
    ellOutline/.style={gDark, line width=0.5pt},
    ellName/.style={font=\sffamily\scriptsize\bfseries, text=gDark, inner sep=1pt},
  ]

% ============================================================
% Outer scope(C) rectangle -- taller now, more breathing room.
% ============================================================
\def\rectXmin{0}
\def\rectXmax{10.20}
\def\rectYmin{0}
\def\rectYmax{7.3}
\def\axisY{-1.30}

\draw[gDark, line width=0.6pt, rounded corners=3pt]
  (\rectXmin, \rectYmin) rectangle (\rectXmax, \rectYmax);

% scope(C) header label, inset inside the top-left corner so the
% rounded rectangle border draws intact behind it.
\node[font=\sffamily\bfseries\footnotesize, anchor=north west, inner sep=4pt,
      text=gDark]
  at (\rectXmin+0.12, \rectYmax-0.12) {$\Scope(C)$};

% ============================================================
% Upper region: anchor-domain and touched-between ellipses.
% Labels sit ABOVE each ellipse (still inside the rectangle).
% ============================================================
\def\anchorCx{3.7}
\def\anchorCy{4.85}
\def\touchedCx{6.5}
\def\touchedCy{4.85}
\def\ellRx{1.95}
\def\ellRy{0.85}

\node[ellName, anchor=south]
  at (\anchorCx-0.65, \anchorCy+\ellRy+0.10) {anchor-domain};
\node[ellName, anchor=south]
  at (\touchedCx+0.65, \touchedCy+\ellRy+0.10) {touched-between};

% Intersection of the two ellipses, filled with gPale to make the
% Venn overlap visually obvious (200 sits in this lens region).
\begin{scope}
  \clip (\anchorCx, \anchorCy) ellipse [x radius=\ellRx, y radius=\ellRy];
  \fill[gPale] (\touchedCx, \touchedCy) ellipse [x radius=\ellRx, y radius=\ellRy];
\end{scope}

\draw[ellOutline]
  (\anchorCx, \anchorCy) ellipse [x radius=\ellRx, y radius=\ellRy];
\draw[ellOutline]
  (\touchedCx, \touchedCy) ellipse [x radius=\ellRx, y radius=\ellRy];

% Keys inside the ellipses
\node[keyDot] at (\anchorCx-1.15, \anchorCy)   {$100$};
\node[keyDot] at ($(\anchorCx, \anchorCy)!0.5!(\touchedCx, \touchedCy)$) {$200$};
\node[keyDot] at (\touchedCx+1.15, \touchedCy) {$300$};

% Union annotation between upper and lower ellipses
\node[font=\sffamily\scriptsize\itshape, text=gMid, anchor=center,
      align=center, inner sep=2pt]
  at (5.10, 3.20)
  {union $=$ table-scope-between $= \{100, 200, 300\}$};

% ============================================================
% Lower region: replay-keys ellipse.
% Label sits ABOVE it (inside the rectangle).
% ============================================================
\def\replayCx{5.10}
\def\replayCy{1.55}
\def\replayRx{3.2}
\def\replayRy{0.70}

\node[ellName, anchor=south]
  at (\replayCx, \replayCy+\replayRy+0.10) {replay-keys};

\draw[ellOutline]
  (\replayCx, \replayCy) ellipse [x radius=\replayRx, y radius=\replayRy];

\node[keyDot] at (\replayCx-2.10, \replayCy) {$100$};
\node[keyDot] at (\replayCx,      \replayCy) {$200$};
\node[keyDot] at (\replayCx+2.10, \replayCy) {$300$};

% ============================================================
% Source-log axis BELOW the rectangle.
% ============================================================
\draw[axisLine] (0.10, \axisY) -- (10.10, \axisY);
\node[sublabelDim, anchor=west, inner sep=0pt]
  at (10.20, \axisY) {source-log};

\def\bX{2.60}
\def\fX{7.60}
\draw[axisTick] (\bX, \axisY+0.07) -- (\bX, \axisY-0.07);
\node[font=\sffamily\scriptsize, anchor=north, inner sep=1pt]
  at (\bX, \axisY-0.07) {$b$};
\draw[axisTick, line width=0.8pt] (\fX, \axisY+0.10) -- (\fX, \axisY-0.10);
\node[font=\sffamily\scriptsize\bfseries, anchor=north, inner sep=1pt]
  at (\fX, \axisY-0.07) {$f = \Frontier(C)$};

% Relation b <= f.  Lifted to y+0.30 so the conjunct-(1) leader (which
% lands at y+0.10) clears it.
\node[font=\sffamily\scriptsize\itshape, text=gMid, anchor=south, inner sep=2pt]
  at ($(\bX, \axisY)!0.5!(\fX, \axisY) + (0, 0.30)$) {$b \le \Frontier(C)$};

% ============================================================
% Conjunct circled numerals.
% ============================================================
% (1) lifted slightly above the axis, with a leader pointing to the
% interval between b and f (the relation, not the point b alone).
\node[conjNum] (c1n) at (\bX-1.40, \axisY+0.45) {1};
\draw[depFlow, line width=0.4pt]
  (c1n.east) to[bend left=12]
  ($(\bX, \axisY)!0.5!(\fX, \axisY) + (0, 0.10)$);

% (2) upper-left, leader into anchor-domain ellipse.
% Pinned at absolute Y=6.15 so the rectYmax bump doesn't push it
% closer to the scope(C) header above.
\node[conjNum] (c2n) at (\rectXmin+0.65, \rectYmax-1.15) {2};
\draw[-{Latex[length=3pt, width=2.5pt]}, gMid, line width=0.4pt]
  (c2n.south east) to[bend right=10] ($(\anchorCx-\ellRx, \anchorCy+0.45)$);

% (3) upper-right, leader into touched-between.  Same pinning as (2).
\node[conjNum] (c3n) at (\rectXmax-0.65, \rectYmax-1.15) {3};
\draw[-{Latex[length=3pt, width=2.5pt]}, gMid, line width=0.4pt]
  (c3n.south west) to[bend left=10] ($(\touchedCx+\ellRx, \touchedCy+0.45)$);

% (4) lower-right, leader into replay-keys
\node[conjNum] (c4n) at (\rectXmax-0.65, \rectYmin+0.83) {4};
\draw[-{Latex[length=3pt, width=2.5pt]}, gMid, line width=0.4pt]
  (c4n.north west) to[bend right=10] ($(\replayCx+\replayRx, \replayCy-0.15)$);

% ============================================================
% Header summarizing the predicate as a 4-conjunct AND.
% Sits ABOVE the rectangle, right-aligned.
% ============================================================
\node[font=\sffamily\scriptsize, text=gDark, anchor=south east, align=right,
      inner sep=2pt]
  at (\rectXmax, \rectYmax+0.10) {%
  \textbf{$\WholeTableClaimScope(C, H, b) \;=\;
          (1) \land (2) \land (3) \land (4)$}};

\end{tikzpicture}

\caption{Anchor domain and whole-table claim scope. The four numbered
conjuncts of $\WholeTableClaimScope(C, H, b)$ are annotated on the
diagram. The base state $b_0$, fixed by context, is implicit in
$\WholeTableClaimScope$ and $\AnchorDomain$. \textbf{(1)} The anchor $b$ precedes the frontier on the
source-log axis: $b \le \Frontier(C)$.
\textbf{(2)} $\AnchorDomain(H, b) \subseteq \Scope(C)$.
\textbf{(3)} $\TouchedBetween(H, b, \Frontier(C)) \subseteq \Scope(C)$.
\textbf{(4)} $\ReplayKeys(\CleanPrefix(C)) \subseteq \Scope(C)$.
Concrete keys follow the worked example of Section~\ref{a-small-anchor-example}: $100$ sits in
anchor-domain only, $200$ in the overlap, $300$ in touched-between
only. Replay-keys are $\{100, 200, 300\}$.}
\label{fig:whole-table-claim-scope}
